% GENERATED FILE. Edit canonical lessons, then run npm run generate:books.
% canonical-source-hash: fnv1a64:641b0d9fc7807d4b
% canonical-lessons: ZH-C462-jingxi, ZH-C462-yixiangbudao, ZH-C462-hen, ZH-C462-sixin, ZH-C462-bianxin

\chapter{Surprises}
\label{ch:zh-surprises}
\hlchaptermodality{462}

\begin{chapteropening}
\textbf{By the end of this chapter:} \emph{I can talk about a lovely surprise, something unexpected, hating someone, giving up hope and falling out of love.}

Complete the last lesson of chapter 462: I can talk about a lovely surprise, something unexpected, hating someone, giving up hope and falling out of love.
\end{chapteropening}

\section[jīngxǐ]{jīngxǐ — pleasant surprise}

\label{lesson:ZH-C462-jingxi}



\begin{quote}
Before the new one: say the Mandarin for treasured, then the Mandarin for to look forward to.
\end{quote}

\subsection*{You'll want to know: jīngxǐ}
\emph{jīngxǐ} — \textquotedblleft{}pleasant surprise\textquotedblright{}.

\begin{quote}
\emph{zhè shì wǒ jīnnián zuì dà de jīngxǐ} — This is the biggest surprise I've had all year.
\end{quote}

Its characters come later; for now it is a word for the ear.

\begin{cousinweb}
\emph{jīngxǐ} is built from \emph{jīng} 'to startle' and \zh{喜} 'joy', both inherited from Middle Chinese, \emph{kjaeng} and \emph{xiX} (Guangyun, in Baxter's transcription), which glosses the first 'fear' and, after an older book, a startled horse, and the second 'joy'.
\end{cousinweb}

\begin{grammarlens}[title={What you've built}]
One more word for saying how it was.
\end{grammarlens}

\subsection*{Your turn}
\begin{itemize}
\raggedright
  \item \emph{Say it:} \emph{jīngxǐ}
  \item \emph{Say it:} \emph{jīngxǐ}, once more
  \item \emph{Say it:} \emph{zhè shì wǒ jīnnián zuì dà de jīngxǐ} — This is the biggest surprise I've had all year.
  \item \emph{Recall:} say the Mandarin for treasured, then the Mandarin for to look forward to, then say \emph{jīngxǐ} again
\end{itemize}

\begin{tcolorbox}[breakable,colback=teal!4,colframe=teal!35!black,title={Before you move on}]
What does jīngxǐ mean? (\textquotedblleft{}Pleasant surprise\textquotedblright{}.) Say it once more.
\end{tcolorbox}
\section[\texorpdfstring{\zh{意想不到} (yìxiǎngbúdào)}{意想不到 (yìxiǎngbúdào)}]{\zh{意想不到} (yìxiǎngbúdào) — unexpected; beyond what anyone imagined}

\label{lesson:ZH-C462-yixiangbudao}



\begin{quote}
Before the new one: say the Mandarin for to look forward to, then the Mandarin for a pleasant surprise.
\end{quote}

\subsection*{You'll want to know: \zh{意想不到}}
\textbf{\zh{意想不到}} — \emph{yìxiǎngbúdào} — \textquotedblleft{}unexpected; beyond what anyone imagined\textquotedblright{}.

\begin{quote}
\textbf{\zh{这真是一个意想不到的好方法}} \emph{zhè zhēn shì yí ge yìxiǎngbúdào de hǎo fāngfǎ} — That really is an unexpectedly good way of doing it.
\end{quote}

\zh{不} is said with a rising tone before the falling tone of dào: bú dào.

\begin{cousinweb}
\textbf{\zh{意想不到}} (\emph{yìxiǎngbúdào}) is built from \zh{意} 'the mind', \zh{想} 'to think', \zh{不} 'not' and \zh{到} 'to reach', all inherited from Middle Chinese, \emph{'iH}, \emph{sjangX}, \emph{pjut} and \emph{tawH} (Guangyun, in Baxter's transcription): what the mind's thinking does not reach.
\end{cousinweb}

\begin{grammarlens}[title={What you've built}]
One more word for saying how it was.
\end{grammarlens}

\subsection*{Your turn}
\begin{itemize}
\raggedright
  \item \emph{Say it:} \emph{yìxiǎngbúdào}
  \item \emph{Say it:} \emph{yìxiǎngbúdào}, once more
  \item \emph{Say it:} \emph{zhè zhēn shì yí ge yìxiǎngbúdào de hǎo fāngfǎ} — That really is an unexpectedly good way of doing it.
  \item \emph{Recall:} say the Mandarin for to look forward to, then the Mandarin for a pleasant surprise, then say \emph{yìxiǎngbúdào} again
\end{itemize}

\begin{tcolorbox}[breakable,colback=teal!4,colframe=teal!35!black,title={Before you move on}]
What does \zh{意想不到} mean? (\textquotedblleft{}Unexpected; beyond what anyone imagined\textquotedblright{}.) Say it once more.
\end{tcolorbox}
\section[hèn]{hèn — to hate deeply; to resent}

\label{lesson:ZH-C462-hen}



\begin{quote}
Before the new one: say the Mandarin for a pleasant surprise, then the Mandarin for beyond what anyone imagined.
\end{quote}

\subsection*{You'll want to know: hèn}
\emph{hèn} — \textquotedblleft{}to hate deeply; to resent\textquotedblright{}.

\begin{quote}
\emph{nǐ hái hèn tā ma} — Do you still hate him?
\end{quote}

Its characters come later; for now it is a word for the ear.

\begin{cousinweb}
\emph{hèn} is inherited from Middle Chinese \emph{honH} (Guangyun, in Baxter's transcription), glossed 'resentment'.
\end{cousinweb}

\begin{grammarlens}[title={What you've built}]
One more word for saying how it was.
\end{grammarlens}

\subsection*{Your turn}
\begin{itemize}
\raggedright
  \item \emph{Say it:} \emph{hèn}
  \item \emph{Say it:} \emph{hèn}, once more
  \item \emph{Say it:} \emph{nǐ hái hèn tā ma} — Do you still hate him?
  \item \emph{Recall:} say the Mandarin for a pleasant surprise, then the Mandarin for beyond what anyone imagined, then say \emph{hèn} again
\end{itemize}

\begin{tcolorbox}[breakable,colback=teal!4,colframe=teal!35!black,title={Before you move on}]
What does hèn mean? (\textquotedblleft{}To hate deeply; to resent\textquotedblright{}.) Say it once more.
\end{tcolorbox}
\section[\texorpdfstring{\zh{死心} (sǐxīn)}{死心 (sǐxīn)}]{\zh{死心} (sǐxīn) — to give up hope}

\label{lesson:ZH-C462-sixin}



\begin{quote}
Before the new one: say the Mandarin for beyond what anyone imagined, then the Mandarin for to hate deeply.
\end{quote}

\subsection*{You'll want to know: \zh{死心}}
\textbf{\zh{死心}} — \emph{sǐxīn} — \textquotedblleft{}to give up hope\textquotedblright{}.

\begin{quote}
\emph{wǒ děng le tā sān nián xiànzài sǐxīn le} — I waited three years for him, and now I've given up.
\end{quote}

\begin{cousinweb}
\textbf{\zh{死心}} (\emph{sǐxīn}) is built from \zh{死} 'to die' and \zh{心} 'heart', both inherited from Middle Chinese, \emph{sijX} and \emph{sim} (Guangyun, in Baxter's transcription), which explains the first, after an older book, as coming to an end: a heart that has stopped hoping.
\end{cousinweb}

\begin{grammarlens}[title={What you've built}]
One more word for saying how it was.
\end{grammarlens}

\subsection*{Your turn}
\begin{itemize}
\raggedright
  \item \emph{Say it:} \emph{sǐxīn}
  \item \emph{Say it:} \emph{sǐxīn}, once more
  \item \emph{Say it:} \emph{wǒ děng le tā sān nián xiànzài sǐxīn le} — I waited three years for him, and now I've given up.
  \item \emph{Recall:} say the Mandarin for beyond what anyone imagined, then the Mandarin for to hate deeply, then say \emph{sǐxīn} again
\end{itemize}

\begin{tcolorbox}[breakable,colback=teal!4,colframe=teal!35!black,title={Before you move on}]
What does \zh{死心} mean? (\textquotedblleft{}To give up hope\textquotedblright{}.) Say it once more.
\end{tcolorbox}
\section[\texorpdfstring{\zh{变心} (biànxīn)}{变心 (biànxīn)}]{\zh{变心} (biànxīn) — to fall out of love; to stop being faithful}

\label{lesson:ZH-C462-bianxin}



\begin{quote}
Before the new one: say the Mandarin for to hate deeply, then the Mandarin for to give up hope.
\end{quote}

\subsection*{You'll want to know: \zh{变心}}
\textbf{\zh{变心}} — \emph{biànxīn} — \textquotedblleft{}to fall out of love; to stop being faithful\textquotedblright{}.

\begin{quote}
\textbf{\zh{他没有变心他还爱你}} \emph{tā méiyǒu biànxīn tā hái ài nǐ} — He hasn't fallen out of love with you; he still loves you.
\end{quote}

\begin{cousinweb}
\textbf{\zh{变心}} (\emph{biànxīn}) is built from \zh{变} 'to change' and \zh{心} 'heart', both inherited from Middle Chinese, \emph{pjenH} and \emph{sim} (Guangyun, in Baxter's transcription), which glosses the first 'to change'.
\end{cousinweb}

\begin{grammarlens}[title={What you've built}]
One more word for saying how it was.
\end{grammarlens}

\subsection*{Your turn}
\begin{itemize}
\raggedright
  \item \emph{Say it:} \emph{biànxīn}
  \item \emph{Say it:} \emph{biànxīn}, once more
  \item \emph{Say it:} \emph{tā méiyǒu biànxīn tā hái ài nǐ} — He hasn't fallen out of love with you; he still loves you.
  \item \emph{Recall:} say the Mandarin for to hate deeply, then the Mandarin for to give up hope, then say \emph{biànxīn} again
\end{itemize}

\begin{tcolorbox}[breakable,colback=teal!4,colframe=teal!35!black,title={Before you move on}]
What does \zh{变心} mean? (\textquotedblleft{}To fall out of love; to stop being faithful\textquotedblright{}.) Say it once more.
\end{tcolorbox}
